% !TEX root = ../main.tex
% ---------------------------------------------------------------------
%  B1 -- Grundlagen der Automatisierung (Skript Kap. 1, Folien)
% ---------------------------------------------------------------------
\vlsec{cB1}{B1}{Grundlagen der Automatisierung}
\label{sec:grundlagen}

\SW{Definition} (DIN 19222) Ausrüsten einer Einrichtung, so dass sie
\emph{ganz oder teilweise ohne Mitwirkung des Menschen} geschieht. Der
Mensch legt fest \emph{was} geschieht (Programm), nicht mehr \emph{wie}.

\SW{Vier Gebiete} \T{Maschinen}- (Presse, Roboter), \T{Anlagen}-
(Kraftwerk, Chemie, mehrere vernetzte Regelkreise), \T{Gebäude}- (Heizung,
Lüftung, Licht), \T{Geräte}automatisierung (Insulinpumpe, Saugroboter).
Modul-Fokus: Maschinen und Anlagen.

\subsection{Automatisierungspyramide}
\begin{ztab}{@{}llX@{}}
Ebene & System & Zeit \\ \midrule
Firma & ERP & Tage--Monate \\
Produktion & MES & Stunden--Schichten \\
Anlage & Prozessleitsystem, SCADA & s--min \\
System & SPS (Station, Achse) & Sekunden \\
Sensor/Aktor & Regelung (PID) & ms \\
\end{ztab}
Nach \emph{unten}: mehr Datenpunkte, schnellere Reaktion. Nach \emph{oben}:
höherer Aggregationsgrad (Tausende Sensorwerte \ra eine Kennzahl).

\begin{formel}[Steuern vs.\ Regeln]
\T{Steuerung} = offener Wirkungsablauf (Steuerkette): Eingang beeinflusst
Ausgang, \emph{keine Rückmeldung}, ob das Ziel erreicht wurde
(Ablaufsteuerung, Zustandsautomat).\\
\T{Regelung} = geschlossener Wirkungskreis: Istwert wird gemessen, mit
Sollwert verglichen, die Regelabweichung wirkt zurück (PID).
\end{formel}
\SW{Beispiel Abkantpresse} Stempel über 2 Hydraulikzylinder geregelt;
Biegehilfen: 3 Achsen $\times$ 3 Kaskadenregler (Strom, Geschwindigkeit,
Position) $=9$ PID. Dazu reine Steuerung: Bedienung, HMI, Rezepturen,
Produktionsdaten, Safety (Lichtvorhang, Zweihand).

\subsection{Aufbau eines Automatisierungssystems}
HMI (Mensch) \ra SPS \ra Feldbus \ra Antriebe, Aktoren, Sensoren.
\begin{ztab}{@{}lX@{}}
mit Echtzeit & Regler, Prozessabläufe \\
ohne Echtzeit & HMI, Produktionsplanung, Archivierung \\
\end{ztab}
\SW{Anforderungen} Echtzeitfähigkeit (BS, Speicher), E/A von
Prozesssignalen, Sicherheit/Zuverlässigkeit (Redundanz), Resistenz gegen
Umwelt (Gehäuse, Schaltschrank).

\SW{Bauformen} Control Panel (mit Bildschirm), Compact (Hutschiene, wie
CX im Labor), 19''-Rack-Industrie-PC.
\T{SPS}: Wenn-Dann-Logik + Zustandsspeicher. \T{CNC}: Bahnregelung,
CAD-Import (IGES/DXF). \T{Industrie-PC}: gehärteter PC + Soft-SPS
(TwinCAT) \ra SPS und HMI in einem.

\SW{Schaltschrank} Hauptschalter \ra Sicherungsautomaten \ra Netzteil
(230/400\,V AC \ra 24\,V DC) \ra SPS \ra E/A-Klemmen \ra
Frequenzumrichter (drehzahlvariable Antriebe) \ra Schütze/Relais (grosse
Lasten, trennen Steuer- von Leistungsebene) \ra Not-Aus-Kreis
(Sicherheitsrelais/Safety-SPS). Schutzart z.\,B.\ IP54.

\SW{I/O-Baum} \T{Online} (Scan der laufenden Hardware) oder \T{Offline}
(Klemmen aus Katalog, ESI-Dateien, nach Stückliste). EPLAN-Kopplung:
TF6360 importiert das Elektroschema.

\subsection{Feldbusse}
\label{sec:feldbus}
Gemeinsame Leitung statt Einzelverdrahtung (zentralisiert).
\begin{ztab}{@{}XX@{}}
Vorteile & Nachteile \\ \midrule
weniger Verkabelung & komplexer (Protokoll, Adressen) \\
einfach erweiterbar & teurer, Bus-Analyzer nötig \\
 & ev.\ längere Reaktionszeit \\
\end{ztab}
\SW{Generationen} klassisch: Profibus (RS-485), CAN/CANopen (Auto,
J1939). Ethernet-basiert: \T{EtherCAT} (Beckhoff), Profinet, EtherNet/IP,
SERCOS III. Gebäude: LON, EIB/KNX.

\begin{formel}[EtherCAT vs.\ TCP/IP]
\T{Gleiches Kabel} (Schicht 1, Cat5/RJ45, 100BASE-TX), aber
\emph{kein} IP und \emph{kein} TCP/UDP im Echtzeitpfad: eigener Ethertype
direkt auf Schicht 2, Verarbeitung on-the-fly, feste projektierte
Teilnehmer-Reihenfolge statt Routing, kein CSMA/CD \ra deterministisch,
$\mu$s bis wenige ms.
\end{formel}
\SW{Topologien} Linie (bis 100\,m zwischen Teilnehmern), Baum/Abzweig,
Stern (Junctions), \T{Ring mit Kabelredundanz} (letzter Teilnehmer zurück
zum Master: Kabelbruch \ra Gegenrichtung übernimmt). Fremdgeräte direkt
mit EtherCAT (Schunk-Greifer, Bota-Kraftsensor).

\subsection{Programmiersprachen IEC 61131-3}
\label{sec:sprachen}
\begin{ztab}{@{}llX@{}}
Sprache & Typ & Merkmal \\ \midrule
AWL / IL & textuell & SPS-Assembler, keine Schleifen, veraltet \\
KOP / LD & grafisch & Ersatz für Relaislogik, nicht für Komplexes \\
FBS / FBD & grafisch & Bit-Verknüpfungen, Datenfluss übersichtlich \\
\T{ST} & textuell & Hochsprache (Pascal/C), Schleifen, in Europa verbreitet \\
AS / SFC & grafisch & Zustandsabläufe, gut lesbar \\
\end{ztab}
\begin{falle}[Anzahl Sprachen]
Es sind \T{fünf} Sprachen \emph{insgesamt} -- ST plus \emph{vier} weitere.
IEC 61131-3 ist die einzige weltweit gültige Norm und
herstellerübergreifend (nicht nur Beckhoff).
\end{falle}
\SW{Warum ST} Filter, PID, Automaten brauchen Schleifen, Arrays, REAL,
Fallunterscheidungen. Reiner Text \ra Diff, Git, Code-Review.
\SW{Ausblick} \T{EN 61499}: objektorientierte Weiterentwicklung, verteilte
ereignisgesteuerte FBs über mehrere Steuerungen. \T{Codegenerierung}:
Simulink PLC Coder, TwinCAT-Simulink, B\&R Automation Studio --
validierter Regler ohne Abschreibfehler, aber weniger Kontrolle über den
Code. Direkter vs.\ indirekter Reglerentwurf siehe \ref{sec:pid}.
